% !TEX root = ../main_en.tex

\section{Introduction}\label{sec:intro}

\subsection{Research Background}

\textbf{Active matter} refers to particulate systems driven out of equilibrium at the particle scale\cite{boltz2025}. Such systems exhibit many collective phenomena that are inaccessible to equilibrium physics, the most famous example being long-range orientational order in two dimensions\cite{vicsek1995}.

In general, activity (particularly the self-propulsion of particles) is a means of establishing long-range correlations in dilute systems, even when interactions are short-ranged. This property holds even in the absence of global spatial order, orientational order, or any other type of order. Taking this to the extreme: even when the nonequilibrium steady state is characterized by a uniform \textbf{single-particle distribution function}, correlations between particles remain crucial for understanding the behavior of the system.

\begin{notebox}[Key Physical Picture]
These correlations quantify the statistical deviation of positional configurations from realizations of a simple \textbf{homogeneous Poisson point process}. Even in the absence of any macroscopic order, long-range correlations in active systems can still lead to significant corrections to statistical fluctuations.
\end{notebox}

\subsection{Hyperuniformity}

\textbf{Hyperuniform configurations}\cite{torquato2018} are a special class of disordered point patterns in which the particle number fluctuations of large subsystems exhibit \textbf{anomalous scaling behavior}:
\begin{equation}\label{eq:hyperuniform-def}
\Delta N \sim r^\beta \quad (r \to \infty), \quad d-1 < \beta < d
\end{equation}
where $d$ is the spatial dimension. This stands in sharp contrast to the bulk law of ordinary Poisson processes, $\Delta N \sim \sqrt{N} \sim r^{d/2}$.

\textbf{Complementary definition}: Hyperuniformity can also be defined via the static structure factor:
\begin{equation}\label{eq:structure-factor-def}
S(\mathbf{k}) = \langle \delta\tilde{\rho}(\mathbf{k})\, \delta\tilde{\rho}(-\mathbf{k}) \rangle
\end{equation}
When $S(\mathbf{k} \to \mathbf{0}) \to 0$, the system is in a hyperuniform state.

\subsection{Antialigning Interactions}

\textbf{Antialigning} orientational interactions are less studied than aligning interactions, but exist in certain experiments\cite{boltz2025}. For dilute systems, density fluctuations cause particles to interact and typically produce opposite orientations. Therefore, \textbf{center-of-mass motion} is suppressed (though not completely forbidden by conservation laws), and successive interactions of the same particle are also suppressed.

\begin{notebox}[Core Intuition]
The core intuition of this paper is: in systems of self-driven particles with short-range antialigning interactions, density fluctuations are ``dissolved'' in a nonlocal manner. Similar to Manna processes or adsorbed-state systems with center-of-mass conservation, the correlation mechanism in antialigning systems can effectively suppress large-scale density fluctuations.
\end{notebox}

\subsection{Main Contributions of the Paper}

The main contributions of this paper include:

\begin{enumerate}
    \item \textbf{One-dimensional analytical model}: A one-dimensional active lattice gas model is constructed and analyzed rigorously using the \textbf{Poisson representation}
    \item \textbf{Fluctuating hydrodynamics}: Fluctuating hydrodynamic equations for the Poisson fields are derived
    \item \textbf{Discovery of imaginary noise}: \textbf{Imaginary noise} describing non-Poissonian statistics is discovered, marking the non-Poissonian character of particle number fluctuations
    \item \textbf{Analytical calculation of the structure factor}: The nontrivial structure factor $S(k)$ is calculated analytically
    \item \textbf{Two-dimensional numerical verification}: Theoretical predictions are verified through simulations of an off-lattice Vicsek-like model
\end{enumerate}

These notes focus on filling in the theoretical derivation details omitted in the original article, particularly:
\begin{itemize}
    \item Construction of the master equation and complete derivation of the Poisson representation
    \item Detailed construction of the Fokker-Planck equation
    \item Derivation process of the Langevin equations
    \item Analytical calculation of the structure factor
\end{itemize}
